\section{Lecture 6: Regularization and Structure}

\subsection{Motivation}

In the previous lecture, we studied generalization and observed that:

\begin{itemize}
    \item Models with high capacity can overfit
    \item Low training error does not guarantee low test error
\end{itemize}

\medskip

\noindent
This raises a central question:

\medskip

\noindent
\textit{How can we guide learning toward solutions that generalize well?}

\medskip

\noindent
The answer lies in \emph{regularization}, which introduces additional structure or constraints into the learning process.

\medskip

\noindent
\textbf{Guiding idea.}  
Regularization shapes the solution space to favor functions with desirable properties.

\subsection{Norm-Based Regularization}

A common approach is to penalize the size of the parameters.

\medskip

\noindent
\textbf{Regularized objective.}
\[
\min_\theta \; L(\theta) + \lambda \|\theta\|_p,
\]
where $\lambda > 0$ controls the strength of regularization.

\medskip

\noindent
\textbf{Examples.}

\begin{itemize}
    \item \textbf{$\ell_2$ regularization (ridge):}
    \[
    \Omega(\theta) = \|\theta\|_2^2
    \]

    \item \textbf{$\ell_1$ regularization (lasso):}
    \[
    \Omega(\theta) = \|\theta\|_1
    \]
\end{itemize}

\medskip

\noindent
\textbf{Effects.}
\begin{itemize}
    \item $\ell_2$: discourages large weights, promotes smooth solutions
    \item $\ell_1$: encourages sparsity
\end{itemize}

\medskip

\noindent
\textbf{Interpretation.}
\begin{itemize}
    \item Penalizing large parameters reduces model complexity
    \item Encourages simpler functions
\end{itemize}

\medskip

\noindent
\textbf{Insight.}  
Norm-based regularization controls capacity by constraining parameter magnitude.

\subsection{Constraints and Penalization}

Regularization can also be formulated as a constrained optimization problem.

\medskip

\noindent
\textbf{Constrained formulation.}
\[
\min_\theta \; L(\theta) \quad \text{subject to} \quad \|\theta\|_p \leq C.
\]

\medskip

\noindent
\textbf{Penalized formulation.}
\[
\min_\theta \; L(\theta) + \lambda \|\theta\|_p.
\]

\medskip

\noindent
\textbf{Equivalence.}
\begin{itemize}
    \item Under suitable conditions, these formulations are equivalent
    \item $\lambda$ corresponds to a constraint level $C$
\end{itemize}

\medskip

\noindent
\textbf{Geometric interpretation.}
\begin{itemize}
    \item Optimization occurs within a constraint set (e.g., a ball)
    \item The solution lies where level sets of $L(\theta)$ touch the constraint boundary
\end{itemize}

\medskip

\noindent
\textbf{Example.}
\begin{itemize}
    \item $\ell_2$ constraint: spherical region
    \item $\ell_1$ constraint: diamond-shaped region (promotes sparsity)
\end{itemize}

\medskip

\noindent
\textbf{Insight.}  
Regularization can be viewed as restricting the feasible set of solutions.

\subsection{Implicit Regularization}

Regularization is not only imposed explicitly; it can also arise implicitly through the learning process.

\medskip

\noindent
\textbf{Definition.}  
Implicit regularization refers to biases introduced by the optimization algorithm or parameterization, even without explicit penalties.

\medskip

\noindent
\textbf{Examples.}

\begin{itemize}
    \item Gradient descent in linear models tends to find minimum-norm solutions
    \item Optimization dynamics favor certain regions of parameter space
\end{itemize}

\medskip

\noindent
\textbf{Overparameterized models.}
\begin{itemize}
    \item Many solutions achieve zero training error
    \item Optimization selects one among them
\end{itemize}

\medskip

\noindent
\textbf{Key phenomenon.}
\begin{itemize}
    \item Even without explicit regularization, learned solutions can generalize well
\end{itemize}

\medskip

\noindent
\textbf{Insight.}  
The optimization algorithm itself acts as a form of inductive bias.

\subsection{Early Stopping and Practical Methods}

Regularization can also be implemented through practical training strategies.

\medskip

\noindent
\textbf{Early stopping.}
\begin{itemize}
    \item Stop training before convergence
    \item Prevents overfitting to noise
\end{itemize}

\medskip

\noindent
\textbf{Interpretation.}
\begin{itemize}
    \item Early iterations capture dominant patterns
    \item Later iterations fit noise
\end{itemize}

\medskip

\noindent
\textbf{Other methods.}
\begin{itemize}
    \item Noise injection
    \item Data augmentation
    \item Dropout (in neural networks)
\end{itemize}

\medskip

\noindent
\textbf{Connection to optimization.}
\begin{itemize}
    \item Training trajectory influences final solution
    \item Regularization is tied to dynamics, not just objectives
\end{itemize}

\medskip

\noindent
\textbf{Insight.}  
Regularization can be achieved through both objective design and training procedures.

\subsection{A Unifying Perspective}

Regularization integrates several aspects of learning:

\begin{itemize}
    \item \textbf{Explicit:} penalties and constraints
    \item \textbf{Implicit:} optimization dynamics
    \item \textbf{Practical:} training strategies
\end{itemize}

\medskip

\noindent
\textbf{Core idea.}  
Regularization shapes the set of solutions considered by the learning algorithm.

\medskip

\noindent
\textbf{Key insight.}  
Generalization depends not only on model capacity, but on how solutions are selected within that capacity.

\subsection{Outlook}

In this lecture, we studied:

\begin{itemize}
    \item norm-based regularization,
    \item constrained and penalized formulations,
    \item implicit regularization,
    \item practical regularization methods.
\end{itemize}

\medskip

\noindent
These ideas provide tools for controlling generalization in learning systems.

\medskip

\noindent
In the next part of the course, we turn to concrete models, beginning with linear regression, where these concepts can be analyzed explicitly.

\medskip

\noindent
\textbf{Guiding principle.}  
Effective learning requires not only fitting data, but shaping the solution toward structure and simplicity.